\section{A conical solver without a pure-carrier restriction}
\label{sec:r11-general-conical-solver}

\paragraph{Status and construction.}
\textbf{PROVED THIS ROUND}, submitted for independent exact-source
review.  The constructor is \texttt{/root/r11\_kernel}.
The new formula permits a nonzero first quadruple component and
therefore removes the pure-carrier assumptions of CE5--CE8.
It retains the exact finite native conical algebra, the two Euler
weights, and the two independently verified curl hypotheses.
It does not construct the missing global higher-rank curl primitives.

\subsection{The supplied finite native algebra}

Use precisely CE1--CE4: finitely many continuous homogeneous
piecewise-linear functions on a connected conical open, one finite
common central refinement, and their finite actual differential
polynomials.  The odd operations $I',I''$ are evaluated by
\[
 I'(d'h)=h,\qquad I''(d''h)=h,
 \qquad I'h=I''h=I'(d''h)=I''(d'h)=I'Dh=I''Dh=0,
 \tag{CGS1}
\]
with the total graded Leibniz rule.  They preserve all native
actual relations by the exact CE carrier proof.  $N$ is the
even prime-to-double-prime exterior derivation.
The identities needed here are
\[
 \begin{gathered}
 d'I'+I'd'=E',\qquad d''I''+I''d''=E'',\qquad
 d''I'+I'd''=N,\qquad [N,d']=d'',\qquad[N,d'']=0.
 \end{gathered}
 \tag{CGS2}
\]
The weight conventions are those of CE3: a homogeneous scalar
has weight $(1,1)$, its prime and double-prime first derivatives
have weights $(1,0)$ and $(0,1)$, and a corner factor has weight
$(0,0)$.  These are not total bidegrees.

Fix $\ell\ge0$ and supply an actual triple
\[
 u=(A,B,M)\in L_2^{\ell+2},\qquad
 \deg(A,B,M)=
 ((\ell,\ell+2),(\ell+2,\ell),(\ell+1,\ell+1)).
 \tag{CGS3}
\]
Assume that all three entries lie in the stated finite algebra,
are simultaneous Euler eigenvectors with weights
\[
                    (1,3),\qquad(3,1),\qquad(2,2),
 \tag{CGS4}
\]
and satisfy the four entire actual identities
\[
 d'A+d''M=0,\qquad d''B+d'M=0,
                     \qquad d''A=0,\qquad d'B=0.
 \tag{CGS5}
\]
No restriction is made on their carrier codimensions or
tangential coefficient degrees beyond their actual bidegrees.

\begin{theorem}[Finite general conical third witness]
Define
\[
\begin{gathered}
 P_0=\frac13 I''A,\qquad Q_0=\frac13 I'B,
 \qquad M_0=M-d'P_0-d''Q_0,\\
 P_1=\frac12 I'M_0,\qquad
 U_1=-\frac12 I'NM_0.
\end{gathered}
 \tag{CGS6}
\]
Then every entry is a finite actual expression, and
\[
 \boxed{\qquad
 u=\partial_3\bigl(U_1,0,P_0+P_1,Q_0\bigr),
 \qquad
 \partial_3(U,V,P,Q)=(d'U+d''P,d''V+d'Q,d'P+d''Q).
 \qquad}
 \tag{CGS7}
\]
The denominators are independent of the dimension, rank and
carrier codimension.
\end{theorem}

\begin{proof}
The first two identities in (CGS2), the last two assumptions in
(CGS5), and the weights give
\[
                   d''P_0=A,\qquad d'Q_0=B.
 \tag{CGS8}
\]
The first two assumptions in (CGS5) then imply, with the full
anticommutation signs,
\[
                          d'M_0=d''M_0=0.
 \tag{CGS9}
\]
Its Euler weights remain $(2,2)$: $I''$ raises the first
weight by one, $I'$ raises the second by one, $d'$ lowers
the second by one, and $d''$ lowers the first by one.

Apply (CGS2) to $M_0$ to obtain
\[
                   d'P_1=M_0,\qquad d''P_1=\frac12 NM_0.
 \tag{CGS10}
\]
The second expression need not be zero.  Instead use
\[
 d'NM_0=Nd'M_0-d''M_0=0,
                       \qquad E'(NM_0)=NM_0.
 \tag{CGS11}
\]
The latter follows because $N$ lowers the prime Euler weight
by one.  Cartan's identity therefore yields
\[
                          d'U_1=-\frac12 NM_0.
 \tag{CGS12}
\]
The first component of (CGS7) is the sum of (CGS12), the first
identity of (CGS8), and the second identity of (CGS10).  The
two $NM_0/2$ terms cancel, leaving $A$.  The second component
is the second identity of (CGS8), and the third is the
definition of $M_0$ together with the first identity of
(CGS10).  This proves all three components.

In the pure-carrier setting of CE, $I'NM_0=0$, and the
additional component $U_1$ vanishes.  Thus the exact CE
formula is recovered, not contradicted.  If $\ell=0$,
$U_1$ has negative prime degree and is automatically zero.
No negative-degree primitive is supplied at that endpoint.
\end{proof}

The new term $U_1$ is essential to the statement without a
pure-carrier restriction.  Simply using the older $U_1=0$
formula after a correction has introduced lower-codimension
terms would leave the nonzero first component $NM_0/2$.

\subsection{What an actual curl correction must supply}

Suppose now that $u$ satisfies the finite-algebra and weight
conditions (CGS3)--(CGS4) and is $\partial_2$-closed, but that
its two curls have not been shown to vanish.  The following
is a precise \textbf{CONDITIONAL} interface.

Supply actual forms in the same finite conical algebra
\[
 \begin{array}{c|c|c}
       &\text{bidegree}&\text{Euler weights}\\\hline
 F&(\ell-1,\ell+2)&(1,4)\\
 G&(\ell+2,\ell-1)&(4,1)
 \end{array}
 \tag{CGS13}
\]
with the two literal full identities
\[
                         d''A=DF,\qquad d'B=DG.
 \tag{CGS14}
\]
Their construction, not merely closedness or a finite-period
test, is required.  Define
\[
 u^\circ=u-\partial_3(-F,G,0,0)
                       =(A+d'F,\ B-d''G,\ M).
 \tag{CGS15}
\]
Since $d''d'=-D$, both curls of $u^\circ$ vanish.  It remains
$\partial_2$-closed, and (CGS13) ensures exactly the weights
(CGS4) after correction.  Apply (CGS6) to $u^\circ$, writing
its outputs $U_1^\circ,P_0^\circ,P_1^\circ,Q_0^\circ$.
Then
\[
 \boxed{\quad
 u=\partial_3\left(
 U_1^\circ-F,\ G,\ P_0^\circ+P_1^\circ,\ Q_0^\circ
                  \right).
 \quad}
 \tag{CGS16}
\]
Unlike the pure-carrier version, this formula permits the
correction $d'F$ or $d''G$ to add lower-codimension terms.

For $\ell=0$ both forms in (CGS13) have a negative component
of their bidegree and are zero; this interface then requires
the original curls themselves to vanish.  For $\ell>0$
there is no assertion that (CGS14) is solvable.  A supplied
native current primitive without a verified finite actual
expression would not satisfy the hypotheses.

For a raw-cube application, one may take
$u=\mathcal Q-R^\ell E_{abc}$ only after placing this
\emph{entire} expression in the supplied finite conical
algebra with the stated weights.  Formula (CGS16) would then
give a relative third witness with the actual extending
triple $R^\ell E_{abc}$.  It is not an existence proof for
the two missing curl primitives or for the required native
substitution interface.

\subsection{The small nonzero curl has a local actual primitive}

There is a concrete useful calibration of the distinction
between nonzero curl and obstructed curl correction.  Use
exactly the rank-four ordinary box $O$ of CF13--CF18 and
the original-$K$ coordinates there.  Put
\[
            u_0=\min(x_1,x_2),\qquad k=b-a,
                        \qquad C_0=-Du_0.
 \tag{CGS17}
\]
On the carrier $H=\{x_1=x_2\}$, write
$dx=d''x_1|_H=d''x_2|_H$.  The explicitly constructed actual
form
\[
 F_{\rm loc}
   =x_1\,d''u_0\wedge d''k\wedge d''c
       -u_0\,d''a\wedge d''b\wedge d''c
                             \in B^{0,3}(O)
 \tag{CGS18}
\]
satisfies
\[
 \begin{split}
 DF_{\rm loc}
  &=C_0\wedge\left[
       dx\wedge(d''b-d''a)\wedge d''c
                     +d''a\wedge d''b\wedge d''c\right]\\
  &=d''\bigl(E_{\rm glob}-C_0E_{abc}\bigr)_A\big|_O.
 \end{split}
 \tag{CGS19}
\]

\begin{proof}
All of $x_1,a,b,c,k$ are original unit valuations, so their
$D$ is zero.  Direct full Leibniz differentiation gives
\[
 D(x_1\,d''u_0\wedge d''k\wedge d''c)
      =-d''x_1\wedge Du_0\wedge d''k\wedge d''c,
 \tag{CGS20}
\]
and
\[
 D(-u_0\,d''a\wedge d''b\wedge d''c)
      =-Du_0\wedge d''a\wedge d''b\wedge d''c.
 \tag{CGS21}
\]
Use $C_0=-Du_0$ and restrict the first coefficient to its
actual carrier $H$.  This yields the first line of (CGS19).
The second line is exactly the full CF21 calculation.
\end{proof}

This primitive is only on the displayed ordinary native box;
it is not a global form on the whole projective angular model.
For example its individual $x_1$ has not been extended across
the corresponding homogeneous-coordinate zero locus.
The correction $A\mapsto A+d'F_{\rm loc}$ removes the first
curl there, but includes lower-codimension terms such as
$d'x_1\wedge d''u_0\wedge d''k\wedge d''c$.  It does not
preserve the pure-carrier hypothesis of the older CE theorem.
The general formula above explains exactly how to accommodate
such terms \emph{if both complete curl corrections are supplied
in one admissible finite conical presentation}.  No second
curl primitive, global continuation, cap overlap gluing or
universal higher-rank certificate is inferred from (CGS18).

\paragraph{Actual operations, refinement and topology.}
The algebra and its equality relation are the native actual
ones of CE, and the final rules in (CGS1)--(CGS16) are finite
scalar, exterior and differential operations.  The CE proof
ensures their compatibility with central refinements and
native actual zero relations.  On each fixed finite whole
presentation they have the corresponding finite coefficient
seminorm bounds.  This is not a topology comparison or a
continuous global section on arbitrary Berkovich forms.

To use a native identity on original analytic data, one must
eliminate every occurrence of $I',I''$ by the finite rules
(CGS1), so that only the supplied actual scalar arguments
and their differentials remain.  Proper-entry zeros and
positive auxiliary caps then require their separately
verified NCP interface, and arbitrary original real cap
constants require the exact same-passive-block CT comparison.
Neither restriction is removed by the algebraic theorem.

\paragraph{Exact dependencies and exclusions.}
The finite conical algebra and Cartan identities are the exact
CE source
\path{68a14ad23c956c211aaae0c188aebb672332181b0a13a269e21cda595a0b6f1d}.
The local calibration uses only CF13--CF21 from
\path{90bc9837cf43de60a2e20c494374b412bb36d68a6d5deb61697fb04bc68d2e52}.
The optional original-field substitution interface is NCP
\path{144df3e33ec6beeefa7907b1292edc31855d289727da777082c5015d33676743},
and the cap comparison is CT
\path{0c10e28336f9c779cc3dc4a35393b1fc52d1f6d7d154a4ec466c13d7c6514e95}.
No MG/FT input, full actual kernel theorem, general matrix-path
classification, field extension, completed form space or
representability from an abstract functional is invoked.
